\documentclass[10pt]{scrartcl}
\usepackage[mdthm,colorsec,noauthor,noasy]{evan}
\usepackage{handout}
\DeclareMathOperator{\im}{range}\DeclareMathOperator{\tr}{tr}
\title{MATH 110 --- Week 9 --- Solutions}
\subtitle{Notes Chapters 16--17; Homework 9}
\author{}\date{}\ihead{}
\begin{document}\maketitle
\begin{prob}\textbf{HW9 P1.}\par
\idea{Construct a basis on which the idempotent is zero or identity.}
\begin{steps}
\ii Every $v=(v-Pv)+Pv$ has first part in $\ker P$, second in $\im P$, because $P(v-Pv)=Pv-P^2v=0$. If $w\in\ker P\cap\im P$, then $w=Pu$ gives $Pw=P^2u=Pu=w$, but also $Pw=0$, so $w=0$. Thus $V=\ker P\oplus\im P$ (also HW6 P1(b)).
\ii Choose bases $u_1,\ldots,u_k$ of $\ker P$ and $w_1,\ldots,w_r$ of $\im P$, and concatenate them. Their images are $Pu_i=0$ and $Pw_j=w_j$, so the matrix is $\operatorname{diag}(0_k,I_r)$.
\ii Its trace is $r1_F$, where $r=\operatorname{rank}P$. Over $\mathbb R$ or $\mathbb C$ this is the usual number $r$; in positive characteristic it is the integer interpreted as a field scalar.
\end{steps}\end{prob}
\begin{prob}\textbf{HW9 P2.}\par\idea{The largest 2-block must have size 2 and every $-1$-block has size 1.}\begin{steps}\ii If $m$ is the algebraic multiplicity of 2, then $m\ge2$; the presence of $-1$ gives $5-m\ge1$. Hence $m=2,3,4$, with pairs $(m,5-m)=(2,3),(3,2),(4,1)$.\ii Proposition 79 gives $\tr T=2m-(5-m)=3m-5$, yielding $1,4,7$.\ii These are realized respectively by $\operatorname{diag}(J_2(2),[-1],[-1],[-1])$, $\operatorname{diag}(J_2(2),[2],[-1],[-1])$, and $\operatorname{diag}(J_2(2),J_2(2),[-1])$. Each has the required minimal polynomial by Week 8's maximal-block rule.\end{steps}\end{prob}
\begin{prob}\textbf{HW9 P3.}\par\idea{Use eigenvalue sums for nilpotence and direct multiplication for a commutator.}\begin{steps}\ii A nonzero-dimensional complex nilpotent operator has only eigenvalue 0 (Week 8), so Proposition 79 gives trace zero. The zero-dimensional case also has trace zero by the empty sum.\ii $T=\operatorname{diag}(1,-1)$ is nonzero with trace 0 and $T^2=I$, hence is not nilpotent.\ii For $A=\left(\begin{smallmatrix}0&1\\0&0\end{smallmatrix}\right)$, $B=\left(\begin{smallmatrix}0&0\\1&0\end{smallmatrix}\right)$, multiplication gives $AB-BA=\operatorname{diag}(1,-1)\ne0$ and trace 0.\end{steps}\end{prob}
\begin{prob}\textbf{HW9 P4.}\par\idea{Bilinearity and symmetry are easy; positive definiteness distinguishes the three formulas.}\begin{steps}\ii The first is a positive weighted coordinate product with weights 1 and 2, so is an inner product by the verified example.\ii The second fails positivity at $e_2$, giving $-1$. The third fails positive definiteness at $(1,-1)\ne0$, giving 0. Neither is an inner product.\end{steps}\end{prob}
\begin{prob}\textbf{HW9 P5.}\par\idea{A nonzero continuous polynomial has a positive squared integral; parity simplifies the computations.}\begin{steps}\ii Linearity and symmetry follow from integration. A nonzero polynomial has only finitely many roots (Week 4), so choose a point in $(-1,1)$ where $p\ne0$. Continuity then gives a positive-length subinterval with $p^2$ bounded below by a positive number, hence $\int_{-1}^1p^2>0$. Thus the formula is an inner product.\ii $\|1\|^2=2$, $\|x\|^2=2/3$, $\|x^2\|^2=2/5$. Therefore the norms are $\sqrt2,\sqrt{2/3},\sqrt{2/5}$.\ii $\langle1,x\rangle=\int_{-1}^1x=0$ and $\langle x,x^2\rangle=\int_{-1}^1x^3=0$, so those pairs are orthogonal. But $\langle1,x^2\rangle=2/3\ne0$, so that pair is not.\end{steps}\end{prob}
\begin{prob}\textbf{HW9 P6.}\par\idea{The sum is the dot product with a vector of three ones.}\begin{steps}\ii Cauchy--Schwarz gives $|a+b+c|\le\sqrt{a^2+b^2+c^2}\sqrt3$. Squaring gives the requested inequality.\ii Equality holds exactly when $(a,b,c)$ and $(1,1,1)$ are dependent, which is exactly $a=b=c$, including all zero.\end{steps}\end{prob}
\begin{prob}\textbf{HW9 P7 (optional).}\par\idea{Adding opposite-sign squared lengths cancels the mixed terms.}\begin{steps}\ii $\|u+v\|^2=\|u\|^2+\langle u,v\rangle+\langle v,u\rangle+\|v\|^2$.\ii $\|u-v\|^2=\|u\|^2-\langle u,v\rangle-\langle v,u\rangle+\|v\|^2$. Add to obtain $2\|u\|^2+2\|v\|^2$.\end{steps}\end{prob}
\begin{prob}\textbf{LADR 8D Exercise 1.}\par\idea{Write the rank-one operator's matrix in an arbitrary basis and sum its diagonal.}\begin{steps}\ii Choose a basis $b_1,\ldots,b_n$ and write $w=\sum_ic_ib_i$. Then $Tb_j=\langle b_j,v\rangle\sum_ic_ib_i$, so its $(i,j)$ entry is $c_i\langle b_j,v\rangle$.\ii Thus $\tr T=\sum_jc_j\langle b_j,v\rangle=\langle\sum_jc_jb_j,v\rangle=\langle w,v\rangle$, by first-position linearity.\end{steps}\end{prob}
\begin{prob}\textbf{LADR 6A Exercise 12.}\par\idea{Pair each positive square root with its reciprocal so the dot product is exactly 4.}\begin{steps}\ii Apply Cauchy--Schwarz to $(\sqrt a,\sqrt b,\sqrt c,\sqrt d)$ and $(1/\sqrt a,1/\sqrt b,1/\sqrt c,1/\sqrt d)$. Their dot product is 4, so squaring gives the bound 16.\ii Equality requires a scalar $t$ with $\sqrt a=t/\sqrt a$, and similarly for the other entries. Thus $a=b=c=d=t>0$. Conversely these equal values give equality.\end{steps}\end{prob}
\end{document}
